\documentclass{report}
\usepackage[utf8]{inputenc}
\usepackage[spanish]{babel}
\usepackage{amsmath}
\usepackage{amssymb}
\usepackage{amsthm}
\usepackage{float}
\usepackage{tikz}
\usetikzlibrary{shapes}

\theoremstyle{definition}
\newtheorem{definicion}{Definición}[chapter]
\theoremstyle{plain}
\newtheorem{teorema}{Teorema}[chapter]
\newtheorem{lema}{Lema}[chapter]
\theoremstyle{remark}
\newtheorem{ejemplo}{Ejemplo}[chapter]

\begin{document}

\chapter{Preliminares}
\label{cap:preliminares}

En este capítulo se presentan las principales definiciones y conceptos que
serán usados en el resto del trabajo, se parte de las nociones básicas de la
teoría de lenguajes y los autómatas finitos, se introducen los elementos de
complejidad computacional necesarios para clasificar los problemas que se
abordan y se presenta el problema de la satisfacibilidad booleana.

\section{Teoría de Lenguajes}
\label{sec:teoria-lenguajes}

En esta sección se presentan los principales conceptos de la teoría de
lenguajes que sirven de base al contenido de los capítulos y secciones
posteriores.

\subsection{Conceptos básicos}

Los conceptos básicos de la teoría de lenguajes formales son: alfabeto,
cadena y lenguaje. Un \textbf{alfabeto}, denotado como $\Sigma$, es un
conjunto finito y no vacío de símbolos; una \textbf{cadena} es una sucesión
finita de símbolos del alfabeto y un \textbf{lenguaje} es un conjunto de
cadenas definido sobre un alfabeto. Por ejemplo, el alfabeto
$\Sigma = \{1, 0\}$ está formado por los símbolos $0$ y $1$, $11$ y $101$
son cadenas sobre el alfabeto $\Sigma$ y un ejemplo de lenguaje es el
conjunto de cadenas de $0$ y $1$ que terminan en $1$.

El símbolo $\varepsilon$ representa la \textbf{cadena vacía}. La
\textbf{longitud} de una cadena $w$, denotada $|w|$, es la cantidad de
símbolos que la componen; en particular $|\varepsilon| = 0$. Dadas dos
cadenas $w_1$ y $w_2$, su \textbf{concatenación} $w_1 w_2$ es la cadena que
se obtiene de escribir los símbolos de $w_1$ seguidos de los símbolos de
$w_2$. Si se tiene una cadena $w$, por convenio $w^n$ representa la cadena
resultante de concatenar $w$ un total de $n$ veces. Si $\Sigma$ es un
alfabeto, $\Sigma^+$ representa a todas las cadenas de uno o más símbolos
que se pueden formar sobre $\Sigma$ y $\Sigma^*$ representa la unión de
$\Sigma^+$ y $\{\varepsilon\}$.

Como los lenguajes son conjuntos, todas las operaciones sobre conjuntos
también se definen para lenguajes: unión, intersección, complemento
[REF]. La siguiente operación, el homomorfismo, permite transformar cadenas
de un alfabeto a otro y es fundamental para relacionar lenguajes definidos
sobre alfabetos distintos.

\begin{definicion}
Dados un alfabeto $\Sigma$ y un alfabeto $\Gamma$, un \textbf{homomorfismo}
es una función:
\[
  h : \Sigma \to \Gamma^*
\]
tal que:
\begin{enumerate}
  \item Para cada $a \in \Sigma$, $h(a)$ es una cadena en $\Gamma^*$, y
  \item si $w = a_1 a_2 \dots a_n$ es una cadena entonces
  \[
    h(w) = h(a_1) h(a_2) \dots h(a_n).
  \]
\end{enumerate}
\end{definicion}

Por ejemplo, si sobre el alfabeto $\Sigma = \{a, b\}$ se define el
homomorfismo $h$, tal que $h(a) = 0$ y $h(b) = 11$, entonces $h(ab) = 011$.

El lenguaje que se obtiene como resultado de aplicarle un homomorfismo $h$
sobre todas las cadenas de un lenguaje $L$ se define como:
\[
  L_h = \{ h(w) \mid w \in L \}.
\]

Por ejemplo, el lenguaje que se obtiene de aplicarle el homomorfismo
anterior al lenguaje
\[
  L = \{ a^n b^n \mid n \in \mathbb{N} \},
\]
es el lenguaje
\[
  L_h = \{ 0^n 1^{2n} \mid n \in \mathbb{N} \}.
\]

Una vez que se tiene definido un lenguaje es posible responder preguntas
como si una cadena dada pertenece al lenguaje, o si el lenguaje es vacío o
contiene algún elemento. Estos problemas se formalizan a continuación.

\begin{definicion}
El \textbf{problema de la palabra} consiste en determinar si una cadena
pertenece a un lenguaje dado.
\end{definicion}

Por ejemplo, dado el lenguaje $L = \{ w \mid \mathrm{last}(w) = 0 \}$,
determinar si $1100100$ pertenece a $L$.

\begin{definicion}
El \textbf{problema del vacío} consiste en determinar si un lenguaje es
vacío.
\end{definicion}

Por ejemplo, determinar si el lenguaje formado por los números pares
mayores que $5$ que son primos es vacío, o si el conjunto formado por
todas las interpretaciones que hacen verdadera a una fórmula booleana es
vacío o no.

En la siguiente sección se definen las gramáticas, un mecanismo que permite
generar los elementos de un lenguaje.

\subsection{Gramáticas}
\label{sec:gramaticas}

Los lenguajes formales pueden describirse mediante \textbf{gramáticas}, un
mecanismo que permite generar los elementos de un lenguaje.

\begin{definicion}
Una \textbf{gramática} es un formalismo utilizado para describir lenguajes
formales. Se define como una $4$-tupla:
\[
  G = (N, \Sigma, P, S),
\]
donde:
\begin{itemize}
  \item $N$: Es un conjunto finito de \textbf{símbolos no terminales}.
  \item $\Sigma$: Es un conjunto finito de \textbf{símbolos terminales},
    que constituyen el alfabeto sobre el que se construyen las cadenas
    del lenguaje. Se cumple que $N \cap \Sigma = \emptyset$.
  \item $P$: Es un conjunto finito de \textbf{reglas de producción} de la
    forma:
    \[
      \alpha \to \beta, \quad \text{donde } \alpha \in (N \cup \Sigma)^*
      \text{ y } \beta \in (N \cup \Sigma)^*.
    \]
  \item $S \in N$: Es el \textbf{símbolo inicial}, que define el punto de
    partida para derivar cadenas del lenguaje.
\end{itemize}
\end{definicion}

Una \textbf{derivación} en la gramática consiste en seleccionar una regla
de producción $\alpha \to \beta$ y sustituir una ocurrencia de $\alpha$
en una cadena $w \in (\Sigma \cup N)^+$ por $\beta$ [REF].

Una cadena $w \in \Sigma^*$ se puede generar por la gramática $G$ si existe
una secuencia de derivaciones que comienza con $S$ y termina con la cadena
$w$. El \textbf{lenguaje generado} por una gramática $G$ se denota como:
\[
  L(G) = \{ w \in \Sigma^* \mid S \xrightarrow{+} w \},
\]
donde $\xrightarrow{+}$ indica una derivación de uno o más pasos.

Un formalismo gramatical también puede definir, para cada cadena de
entrada, si existe una secuencia de derivaciones que permite obtenerla. En
ese caso se dice que la cadena es \textbf{reconocida} por el formalismo, y
el conjunto de todas las cadenas reconocidas conforma el \textbf{lenguaje
reconocido}.

A continuación se presenta la jerarquía de Chomsky, que clasifica a los
lenguajes formales de acuerdo con su poder de generación.

\subsection{Jerarquía de Chomsky}

La \textbf{Jerarquía de Chomsky} [REF] clasifica los lenguajes en cuatro
tipos, según las restricciones en sus reglas de producción y la capacidad
expresiva de los lenguajes que generan.

El primer tipo de gramáticas son las \textbf{Gramáticas irrestrictas}. Estas
gramáticas no tienen restricciones en las reglas de producción. Cada regla
tiene la forma: $\alpha \to \beta$, donde $\alpha, \beta \in (N \cup \Sigma)^*$
y $\alpha \neq \varepsilon$. Todo lenguaje generado por una gramática irrestricta
se denomina \textbf{lenguaje recursivamente enumerable}.

El siguiente tipo en la jerarquía de Chomsky son las \textbf{Gramáticas
dependientes del contexto}. En estas cada regla tiene la forma:
$\alpha A \gamma \to \alpha \beta \gamma$, donde $A \in N$,
$\alpha, \beta, \gamma \in (N \cup \Sigma)^*$, y $|\beta| \ge 1$. Todo lenguaje
generado por una gramática dependiente del contexto se denomina
\textbf{lenguaje dependiente del contexto}. Todo lenguaje dependiente del
contexto es también un lenguaje recursivamente enumerable.

El lenguaje de todas las fórmulas booleanas satisfacibles necesariamente es un
lenguaje dependiente del contexto o un lenguaje recursivamente enumerable [REF].

A continuación le siguen las \textbf{Gramáticas libres del contexto}
(\textit{CFG}). En estas cada regla tiene la forma: $A \to \beta$, donde
$A \in N$ y $\beta \in (N \cup \Sigma)^*$. Todo lenguaje generado por una
gramática libre del contexto se denomina \textbf{lenguaje libre del contexto}.
Todo lenguaje libre del contexto es también un lenguaje dependiente del contexto.
Este tipo de lenguajes permite describir muchas de las propiedades que
generalizan las gramáticas de concatenación de rango [REF]. Además, en [REF]
se usa una CFG para resolver instancias del problema de la satisfacibilidad
booleana.

El último tipo en la jerarquía son las \textbf{Gramáticas regulares}. En estas
las reglas de producción tienen la forma:
\[
  A \to aB \quad \text{o} \quad A \to a,
\]
donde $A, B \in N$ y $a \in \Sigma$. Todo lenguaje generado por una gramática
regular se denomina \textbf{lenguaje regular}. Todo lenguaje regular es un
lenguaje libre del contexto. Esta categoría de lenguajes está relacionada con
los autómatas finitos, los cuales se usan para definir un transductor finito
empleado en este trabajo.

Un lenguaje $A$ es más expresivo que un lenguaje $B$ si todas las cadenas que
pertenecen a $B$ también pertenecen a $A$ y existen cadenas que pertenecen a
$A$ que no pertenecen a $B$; por ejemplo, los lenguajes libres del contexto son
más expresivos que los lenguajes regulares [REF].

La diferencia entre los lenguajes de la jerarquía de Chomsky se puede ilustrar
con el lenguaje \textit{Copy} sobre un alfabeto $\Sigma$, que se define como
$L_{\mathrm{copy}} = \{w^+ \mid w \in \Sigma^*\}$.

Si se toma un caso particular de $L_{\mathrm{copy}}$, al cual se le llama
$L_{\mathrm{copy}}^n = \{w^n \mid w \in \Sigma^*\}$, se cumple que
$L_{\mathrm{copy}}^1$ es un lenguaje regular, mientras $L_{\mathrm{copy}}^k$
$\forall k \ge 2$ es un lenguaje dependiente del contexto [REF].

\begin{figure}[H]
\centering
\begin{tikzpicture}[font=\small]
  \node[above, ellipse, minimum height=2em, minimum width=6em,  draw] (a) {Regulares};
  \node[above, ellipse, minimum height=4em, minimum width=12em, draw] (b) {};
  \node[above, ellipse, minimum height=6em, minimum width=18em, draw] (c) {};
  \node[above, ellipse, minimum height=8em, minimum width=24em, draw] (d) {};
  \path (a.north) node[above] {Libres del contexto}
        (b.north) node[above] {Dependientes del contexto}
        (c.north) node[above] {Recursivamente enumerables};
\end{tikzpicture}
\caption{Jerarquía de Chomsky.}
\label{fig:chomsky}
\end{figure}

En la próxima sección se presentan los autómatas asociados a cada elemento de
la jerarquía de Chomsky.

\section{Autómatas y transductores finitos}
\label{sec:automatas}

Un \textbf{autómata} es una máquina abstracta que procesa cadenas de
símbolos de un alfabeto finito y determina si una cadena pertenece a un
lenguaje [REF].

Cada gramática de la jerarquía de Chomsky tiene un autómata equivalente: los
lenguajes recursivamente enumerables se reconocen por una Máquina de Turing
[REF], los lenguajes dependientes del contexto se reconocen por una Máquina de
Turing linealmente acotada [REF], los lenguajes libres del contexto se
reconocen por un Autómata de pila [REF] y los lenguajes regulares se reconocen
por un Autómata regular [REF].

A continuación se presentan los autómatas finitos y los transductores finitos.

\subsection{Autómata finito}

\begin{definicion}
Un \textbf{autómata finito} [REF] es un modelo matemático que
permite reconocer si una cadena pertenece a un lenguaje regular y se
define como una $5$-tupla
\[
  \mathcal{A} = (Q, \Sigma, \delta, q_0, F),
\]
donde:
\begin{itemize}
  \item $Q$: Es un conjunto finito de \textbf{estados}.
  \item $\Sigma$: Es el \textbf{alfabeto} finito de entrada.
  \item $\delta$: Es la \textbf{función de transición},
    $\delta : Q \times \Sigma \to Q$, que define cómo el autómata cambia
    de estado en función del símbolo leído.
  \item $q_0 \in Q$: Es el \textbf{estado inicial} desde donde comienza la
    computación.
  \item $F \subseteq Q$: Es el conjunto de \textbf{estados de aceptación}
    o \textbf{estados finales}.
\end{itemize}
\end{definicion}

El autómata comienza en el estado inicial $q_0$ y procede a leer el primer
símbolo de la cadena. En cada paso, la función de transición $\delta$
determina, a partir del símbolo actual de la cadena, el siguiente estado
al que debe pasar el autómata. Si al finalizar la lectura del último
símbolo de la cadena, el autómata se encuentra en un estado de aceptación
$q \in F$, entonces la cadena se acepta; en caso contrario, se rechaza.

A continuación se define el transductor finito, una extensión del autómata
finito que produce una cadena de salida durante el reconocimiento.

\subsection{Transductor finito}
\label{sec:transductor}

Un \textbf{transductor finito} [REF] es un modelo computacional
que extiende los autómatas finitos al incluir una salida para cada
transición. Esto permite que al mismo tiempo que se reconoce una cadena
se obtiene otra, que se conoce como \textbf{transducción} de la primera.

\begin{definicion}
Un \textbf{transductor finito} es un autómata finito con una función de
transición extendida que recibe un símbolo de la cadena de entrada y un
estado, y devuelve el estado al cual pasa el transductor y un símbolo.
Como resultado de la aceptación de una cadena el transductor genera una
cadena de salida formada por todos los símbolos que se obtuvieron como
resultado de la función de transición en el proceso de reconocimiento.
Un transductor finito se define como una $6$-tupla:
\[
  T = (Q, \Sigma, \Gamma, \delta, q_0, F),
\]
donde:
\begin{itemize}
  \item $Q$ es el conjunto finito de \textbf{estados}.
  \item $\Sigma$ es el \textbf{alfabeto} de entrada.
  \item $\Gamma$ es el \textbf{alfabeto} de salida.
  \item $\delta : Q \times \Sigma \to Q \times \Gamma^*$ es la
    \textbf{función de transición}, que le asigna a una combinación de
    estado actual y símbolo de entrada un nuevo estado y un símbolo de
    salida.
  \item $q_0 \in Q$ es el \textbf{estado inicial}.
  \item $F \subseteq Q$ es el \textbf{conjunto de estados finales}.
\end{itemize}
\end{definicion}

A modo de ejemplo, se puede definir un transductor finito $T$, que
reconozca cadenas del alfabeto $\{a, b\}$ que tengan ninguna o más $a$
seguida de una o más $b$ y genere una cadena formada por una cantidad de
$0$ igual a la cantidad de $a$ de la cadena original seguida de una
cantidad de $1$ igual a la cantidad de $b$ de la cadena original. Ese
transductor se describe de la siguiente forma:
\[
  T = (Q, \Sigma, \Gamma, \delta, q_0, F),
\]
donde:
\begin{itemize}
  \item $Q = \{q_0, q_1, q_2\}$ es el conjunto de estados.
  \item $\Sigma = \{a, b\}$ es el alfabeto de entrada.
  \item $\Gamma = \{0, 1\}$ es el alfabeto de salida.
  \item $\delta$ es la función de transición definida por:
    \[
      \begin{array}{c|cc}
        \delta & a & b \\ \hline
        q_0 & (q_0, 0) & (q_1, 1) \\
        q_1 & (q_2, 0) & (q_1, 1) \\
        q_2 & (q_2, 0) & (q_2, 1)
      \end{array}
    \]
    Las columnas de la tabla representan el símbolo que se lee y las
    filas el estado en que se encuentra el transductor; los valores de
    la tabla son una tupla que contiene el estado al cual pasa el
    transductor y el símbolo que se escribe.
  \item $q_0$ es el estado inicial.
  \item $F = \{q_1\}$ es el conjunto de estados finales.
\end{itemize}

Para reconocer $aaaabbbb$, primero $T$ empieza por el estado $q_0$ y el
primer carácter $a$, entonces pasa al estado $q_0$ y genera un $0$. Luego,
al leer el segundo carácter $a$, se mantiene en el estado $q_0$ y genera
otro $0$. Este proceso se repite hasta el quinto carácter que es $b$,
entonces pasa al estado $q_1$. Luego, al leer el sexto carácter $b$, se
mantiene en el estado $q_1$ y genera un $1$. Este proceso se repite
varias veces hasta que se llega al final de la cadena, por tanto se
genera la cadena $000011111$.

El lenguaje que se obtiene como resultado de aplicarle un transductor
finito $T$, sobre todas las cadenas de un lenguaje $L$, se define como el
conjunto de todas las cadenas $w$, tales que existe una cadena $e$ para
la cual se cumple que $w$ pertenece al conjunto de cadenas que genera $T$
con la cadena $e$:
\[
  L_T = \{ w \mid \exists e : w \in T(e) \text{ y } e \in L \}.
\]

\section{Complejidad computacional}
\label{sec:complejidad}

En esta sección se definen los principales conceptos de complejidad
computacional: notación asintótica, clases de problemas y reducción
polinomial. Esta definición es importante para analizar la dificultad del
problema de la palabra para los formalismos que se estudian en los
capítulos siguientes.

\subsection{Notación asintótica}

La notación asintótica se utiliza para describir el comportamiento de una
función $f(n)$ a medida que $n$ crece hacia el infinito.

\begin{definicion}
Una función $g(n)$ pertenece a $O(f(n))$ si existen constantes positivas
$c$ y $n_0$ tales que:
\[
  g(n) \le c \cdot f(n) \quad \text{para todo } n \ge n_0.
\]
\end{definicion}

Esta notación proporciona un límite superior asintótico para $g(n)$. Se
dice que un algoritmo tiene un \textbf{tiempo polinomial} si puede
ejecutarse en una complejidad de $O(n^k)$, donde $n$ es el tamaño de la
entrada del algoritmo y $k$ es una constante. La \textbf{complejidad
computacional} de un problema se define como la complejidad del algoritmo
más eficiente que lo resuelve en el peor caso.

\subsection{Clases de problemas}

Los problemas computacionales [REF] se agrupan en diferentes
clases según los recursos necesarios para resolverlos. En este trabajo se
emplean las clases $P$, $NP$, $NP$-Completo y $NP$-Duro.

\begin{definicion}
Un problema pertenece a la clase $\boldsymbol{P}$ si puede resolverse en
tiempo polinomial [REF].
\end{definicion}

\begin{definicion}
Un problema pertenece a la clase $\boldsymbol{NP}$ si su solución puede
verificarse en tiempo polinomial [REF].
\end{definicion}

\begin{definicion}
Un problema pertenece a la clase $\boldsymbol{NP}$\textbf{-Completo}, si
pertenece a $NP$ y además es tan difícil como cualquier otro problema en
$NP$. Esto significa que cualquier problema en $NP$ puede reducirse a
este problema en tiempo polinomial [REF].
\end{definicion}

\begin{definicion}
Un problema pertenece a la clase $\boldsymbol{NP}$\textbf{-Duro}, si es
tan difícil como cualquier otro problema en $NP$, pero no necesariamente
pertenece a $NP$ [REF].
\end{definicion}

La relación entre las clases $P$ y $NP$ es uno de los problemas abiertos
más importantes en la teoría de la computación [REF]. Hasta la
fecha, se desconoce si $P = NP$ o si $P \neq NP$, es decir, no se conoce
si realmente los problemas en $NP$ son más difíciles que los problemas
en $P$. Por otro lado, el conjunto de problemas $NP$-Completo brinda una
base sólida para el problema anterior, ya que dada su definición,
cualquier problema perteneciente a este conjunto que sea soluble en
tiempo polinomial implica que todos los problemas en $NP$ lo son.
Mientras que los problemas en $NP$-Duro pueden resultar aún más difíciles.

Seguidamente se define la reducción polinomial y se precisa el problema
de reconocimiento.

\subsection{Reducción polinomial y problema de reconocimiento}

Una \textbf{reducción polinomial} de un problema $A$ a un problema $B$
es una función computable en tiempo polinomial que transforma instancias
de $A$ en instancias de $B$, tal que la respuesta para la instancia
transformada coincide con la respuesta para la instancia original. Las
reducciones polinomiales permiten transferir resultados de complejidad
entre problemas y son la base para definir la clase $NP$-Completo. En
particular, si un problema $A$ se reduce polinomialmente a $B$ y $B$
pertenece a $P$, entonces $A$ también pertenece a $P$.

La complejidad del \textbf{problema de la palabra} o \textbf{problema de
reconocimiento} para un formalismo gramatical se define como la
complejidad del algoritmo más eficiente que determina, dada una cadena
$w$, si $w$ pertenece al lenguaje que describe dicho formalismo.

Todo problema en Ciencia de la Computación se puede reducir a un problema
de la palabra, ya que cualquier problema se puede codificar como un
lenguaje formal [REF]. Este planteamiento es central en este trabajo:
se utiliza para demostrar que ciertos formalismos gramaticales reconocen
todos los problemas de la clase $NP$-Completo.

\section{Problema de la satisfacibilidad booleana}
\label{sec:sat}

El problema de la satisfacibilidad booleana ($SAT$), es un problema
fundamental en la teoría de la computación y la lógica matemática
[REF]. El objetivo es determinar si existe una asignación de
valores a las variables de una fórmula booleana tal que la expresión sea
verdadera.

A continuación se presentan los principales elementos del $SAT$:
variables, literales, cláusula, fórmulas en forma normal conjuntiva y
fórmulas booleanas equivalentes.

\begin{itemize}
  \item \textbf{Variables booleanas:} Una variable booleana es una
    variable que puede tomar uno de dos valores posibles: \textit{true}
    (verdadero) o \textit{false} (falso). Estas variables se utilizan
    para construir expresiones lógicas.
  \item \textbf{Literales:} Un literal es una variable booleana o su
    negación. Formalmente, si $x$ es una variable booleana, entonces $x$
    y $\neg x$ (la negación de $x$) son literales. Un literal puede
    tomar los valores \textit{true} o \textit{false} dependiendo de la
    asignación de valores a las variables.
  \item \textbf{Cláusulas:} Una cláusula es una disyunción (operador
    \textbf{OR}) de uno o más literales. Por ejemplo, la cláusula
    $(x \vee \neg y \vee z)$ es una disyunción de tres literales: $x$,
    $\neg y$ y $z$. Una cláusula es verdadera si al menos uno de sus
    literales es verdadero. Si todos los literales son falsos, la
    cláusula será falsa.
  \item \textbf{Fórmulas en forma normal conjuntiva:} Una fórmula
    booleana en forma normal conjuntiva ($CNF$) es una conjunción
    (operador \textbf{AND}) de cláusulas. En otras palabras, es una
    expresión booleana que se puede escribir como una serie de
    cláusulas unidas por el operador \textbf{AND}. Por ejemplo:
    \[
      (x \vee \neg y \vee z) \wedge (\neg x \vee y) \wedge (x \vee \neg z).
    \]
  \item \textbf{Fórmulas booleanas equivalentes:} Dos fórmulas booleanas
    se consideran equivalentes si, para cualquier asignación de valores
    a sus variables, ambas producen el mismo resultado lógico. Por
    ejemplo, las fórmulas $x \vee (y \wedge z)$ y
    $(x \vee y) \wedge (x \vee z)$ son equivalentes, ya que para
    cualquier combinación de valores $x, y, z$, ambas tienen el mismo
    valor lógico.
\end{itemize}

Una \textbf{interpretación} de una fórmula booleana es una asignación de
valores \textit{true} o \textit{false} a cada una de sus variables. Para
cualquier fórmula booleana existe una fórmula booleana equivalente en
$CNF$ [REF] y el algoritmo para encontrarla es polinomial, por
lo tanto se puede asumir que toda fórmula booleana está en $CNF$.

\begin{definicion}
El problema de la \textbf{satisfacibilidad booleana}, o $SAT$, consiste
en determinar si existe una asignación de valores \textit{true} o
\textit{false} a las variables de una fórmula booleana, tal que la
fórmula completa sea verdadera.
\end{definicion}

Una fórmula booleana en $CNF$ es \textbf{satisfacible} si existe una
asignación de valores a las variables tal que todas las cláusulas de la
fórmula sean verdaderas simultáneamente.

\subsection{$SAT$ como problema $NP$-Completo y variantes polinomiales}
\label{sec:sat-variantes}

El $SAT$ es el primer problema demostrado como $NP$-Completo
[REF] y juega un rol central en la teoría de la complejidad
computacional. Se define en la clase $NP$ porque dada una asignación de
valores a las variables de la fórmula booleana, se puede verificar en
tiempo polinomial si dicha asignación satisface la fórmula.

Además, la prueba de que $SAT$ es $NP$-Completo fue una de las
contribuciones principales de Stephen Cook en 1971 [REF],
marcando el inicio de la teoría de la $NP$-Completitud.

Un $SAT$ con exactamente $n$ variables distintas en cada cláusula se
denomina $n$-$SAT$. Para el problema 2-$SAT$ existe una solución
polinomial que determina si la fórmula booleana es satisfacible o no
[REF], pero para el problema 3-$SAT$ no se conoce ningún
algoritmo polinomial que permita determinar si una fórmula booleana es
satisfacible o no [REF]. Cualquier fórmula booleana del
problema $n$-$SAT$ se puede reducir a una fórmula booleana equivalente
del problema 3-$SAT$, por lo tanto, $SAT$ es equivalente a 3-$SAT$ en
términos de complejidad computacional [REF].

Aunque no se conoce ningún algoritmo polinomial para resolver el problema
$SAT$ en general, existen casos particulares del problema que sí pueden
ser resueltos en tiempo polinomial como el 2-$SAT$, Horn-$SAT$ y
XOR-$SAT$.

El problema \textbf{2-SAT} es una instancia de $SAT$ donde cada
cláusula contiene exactamente dos literales. Este problema puede ser
resuelto en tiempo polinomial mediante una modelación basada en grafos,
utilizando algoritmos como la detección de componentes fuertemente
conexas en el grafo de implicación [REF]. Esta variante es la
que se utiliza como referencia en los capítulos siguientes para analizar
subclases polinomiales del $SAT$ descritas mediante formalismos de la
teoría de lenguajes.

El problema \textbf{Horn-SAT} es una instancia de $SAT$, donde cada
cláusula tiene a lo sumo un literal positivo. Este problema puede ser
resuelto en tiempo polinomial mediante el algoritmo de resolución de
Horn [REF]. El problema \textbf{XOR-SAT} es una instancia
de $SAT$ donde cada cláusula representa una operación $XOR$ sobre los
literales; puede ser resuelto en tiempo polinomial transformando el
problema en un sistema de ecuaciones lineales modulares y aplicando
eliminación de Gauss [REF].

En el próximo capítulo se presentan las gramáticas de concatenación de
rango, se analizan sus propiedades y limitaciones, y se revisan los
antecedentes directos que abordan el $SAT$ mediante formalismos de la
teoría de lenguajes.

\end{document}
